\chapter{Getting grammaticality wrong}
\section{Relative \textit{whose}}

Old English didn't have relative pronouns the way we have them. It just used \textit{þe} to mark subordinate clauses of various kinds, including relative clauses. 

\ea \label{ex:kyninga}
    \gll \textit{þonne} \textit{ealra} \textit{oþra} \textit{kyninga} \textit{þe} \textit{in} \textit{middangearde} \textit{æfre} \textit{wæron} \dots\\
        Then all other kings that in Middle-Earth ever were \dots\\
    \glt `Then all the kings who were ever on earth \dots'
\z

There was \textit{hwa} `who', but it was mainly interrogative. It had a special construction meaning roughly `whoever', as in \textit{whoever enters remains}, but the \textit{þe} in (\ref{ex:kyninga}) couldn't have been replaced with \textit{hwa}. It was only around 1325, in the midst of Middle English, when we start to see a real relative \textit{who} in the way we know it today.

\ea \label{ex:doȝter-wo}
    \gll \textit{He} \textit{nadde} \textit{bote} \textit{an} \textit{doȝter} \textit{wo} \textit{miȝte} \textit{is} \textit{eir} \textit{be}.\\
        He hadn't but one daughter who might his heir be.\\
    \glt `He had no one but a daughter who could be his heir.'
\z

Similarly, there wasn't a relative \textit{whose} in old English. Instead you get constructions like \textit{þe his} `that his' in (\ref{ex:þe-his}).

\ea \label{ex:þe-his}
    \gll \textit{Eaðig} \textit{bið} \textit{se} \textit{wer}, \textit{þe} \textit{his} \textit{tohopa} \textit{bið} \textit{to} \textit{Drihtne}\\
        Blessed be the man that his hope is in Lord.\\
    \glt `Blessed be the man whose hope is in the Lord'
\z